%%
%% fall-lecture01.tex
%% 
%% Made by Alex Nelson <pqnelson@gmail.com>
%% Login   <alex@lisp>
%% 
%% Started on  2026-10-01T08:50:05-0700
%% Last update 2026-10-01T08:50:05-0700
%% 

\lecture{}

\begin{node}
Jane Bae is the professor. There is no required text, but there are
books on reserve in the library. The syllabus listed several dozen
books as references, most I have. Two which I do not have: Panton's
book on incompressible flows~\cite{panton2024incompressible}
(specifically the 1995 edition), and Frank White's book on viscous flows~\cite{white1991viscous}.
\end{node}

\begin{node}[Overview]
\begin{enumerate}
\item Difference between slids and liquids.
\item Two different views of fluids (microscopic versus
  continuous). Most of this class will focus on the continuum.
\item Dilute gas limit (physical length scales are large enough for continuum?)
\item Criterion for fluid being homogeneous enough (``has enough
  statistical continuity'')
\end{enumerate}
\end{node}

\begin{node}
Broadly speaking, we all learn there are three\footnote{There are more
phases of matter, but we will not consider situations where plasmas
occur.} phases of matter: solids, liquids, and gases. Liquids and
gases behave similarly in many regards, so we group them together
under the term ``fluid''. So what's the difference between a solid and
a fluid?

Any small change in force exerted on a solid---any stress results in a
strain---there will be a deformationj. The shear stress $\tau$ is
proportional to strain $\gamma$ in Hookean solids
\begin{equation}
\tau\sim\gamma.
\end{equation}
In contrast, a fluid has the strain change over time:
\begin{equation}
\tau\sim\frac{\D\gamma}{\D t}.
\end{equation}
There are some exceptions (e.g., with bulk forces).
\end{node}

\begin{node}
What's the difference between liquids and gases? There are three main
differences:
\begin{enumerate}
\item Gases are much more compressible. Compressibility $\kappa$ is defined as
  directly proportional to the partial derivative of density $\rho$
  with respect to pressure $p$ at constant temperature $T$:
  \begin{equation}
\kappa=\frac{1}{\rho}\left(\frac{\partial\rho}{\partial p}\right)_{T}.
  \end{equation}
  For example, $\kappa_{\text{water}}\sim10^{-5}~\mathrm{atm}^{-1}$
  whereas $\kappa_{\text{air}}\sim1~\mathrm{atm}^{-1}$.
\item Liquids form a free surface, whereas gases tend to be space-filling.
\item Liquids are a ``partially ordered system'', the molecular
  spacing $\sigma_{\text{liquid}}\sim\bigOh(1d)$ molecular diameters $d$
  in a liquid. A gas has very little correlation, $\sigma_{\text{gas}}\gtrsim\bigOh(10d)$
\end{enumerate}
Because of this, it is easier to develop a molecular model of a gas,
which fortunately generalizes to liquids.
\end{node}

\subsection{Molecular models of gases}

\begin{node}
The continuum approach: we have macroscopic quantities like velocity,
density, pressure, etc. These are treated as fields. Velocity is a
time-dependent vector field $\vec{v}(\vec{x},t)$. Similarly,
density is a time-dependent scalar field $\rho(\vec{x},t)$, and
pressure is a time-dependent scalar field $p(\vec{x},t)$.
\end{node}

\begin{node}
The microscopic approach has individual molecules. Each molecule has
its positive and velocity. We have the velocity vectors $\vec{v}_{1}$,
$\vec{v}_{2}$, \dots, $\vec{v}_{N}$ for $N$ particles. When $N\gg1$,
we can give up and try treat it as a continuum.

Well, there are three questions we must ask when we want to try to use
the continuum approach:
\begin{enumerate}
\item Is the fluid dilute enough?
\item Are the physical length scales large enough?
\item Is the fluid homogeneous?
\end{enumerate}
The answer to all three questions must be ``Yes'' for us to use the
continuum approach.
\end{node}

\subsection*{Dilute gas limit}

\begin{node}[Dilute gas limit]
Recall Avogadro's law: the volume occupied by a mole of gas at fixed
temperature $T$ and pressure $p$ is constant. We denote the
\define{Number Density} by
\begin{equation}
n:=\frac{(\mbox{number of molecules})}{(\mbox{unit volume})}.
\end{equation}
We can use this to write the mass density $\rho$ as
\begin{equation}
\rho = m\cdot n
\end{equation}
where $m$ is the molecular mass.
\end{node}

\begin{definition}
The \define{Mean Molecular Spacing} is the length $\delta := n^{-1/3}$.
Since
\begin{equation}
n^{-1}=\frac{\mbox{(unit volume)}}{\mbox{(number of molecules)}}\approx\mbox{volume per molecule},
\end{equation}
we see that $n^{-1/3}$ gives a length scale on the order of magnitude
of the space between molecules.
\end{definition}

\begin{definition}
We define the \define{Molecular size} is the term used by engineers to
refer to the ``kinetic diameter'' $d$ of a molecule, which is the size
of a sphere which leads to a scattering between molecules.
\end{definition}

\begin{definition}
A \define{Dilute Gas} is defined as a substance where $\delta\gg d$
the mean molecular spacing is much greater than the molecular size.
Older textbooks use
\begin{equation}
\frac{\delta}{d}>7,
\end{equation}
but this is a heuristic (other textbooks use 8 or 6). If this fails,
then the molecules are ``correlated'' (i.e., too densely packed together).
\end{definition}

\begin{example}
  Air has
  \begin{equation}
d\approx3.7\times10^{-10}~\mathrm{m},
  \end{equation}
  and at STP
\begin{equation}
n_{0}\approx2.7\times10^{25}~\mathrm{m}^{-3},
\end{equation}
so the mean molecular spacing at STP is
\begin{equation}
\delta_{0}\approx3.3\times10^{-9}~\mathrm{m}.
\end{equation}
Then we can calculate:
\begin{subequations}
  \begin{align}
\frac{\delta}{d}
&=\frac{1}{n^{1/3}d}\\
&=\left(\frac{n_{0}}{n}\right)^{1/3}\left(\frac{1}{n_{0}^{1/3}d}\right)\\
&\approx\left(\frac{n_{0}}{n}\right)^{1/3}\cdot9.
  \end{align}
\end{subequations}
Then air is a dilute gas when $\delta/d\gg1$, which is to say
\begin{equation}
\frac{n}{n_{0}}>2.12
\end{equation}
or so.
\end{example}

\subsection*{Physical length scale is large enough}

\begin{definition}
We define the \define{Mean Free Path} $\lambda$ which is the average
distance between collisions (``how far can a molecule go before
hitting another molecule?'').
\end{definition}

\begin{node}[Derivation of mean free path]
If we imagine molecules as spheres with radius $d/2$ half the kinetic
diameter, then they collide if any other sphere [molecule] is within a
radius $d$. We can schematically draw this in two dimensions:
\begin{center}
\includegraphics{img/img.0}
\end{center}
If any molecule's center comes within the dashed circle, then a
collision must occur. What happens in three dimensions? Well, then we
work with the cross-section area $\sigma$ of a sphere of radius $d$, which is
the same as the area of a circle of radius $d$,
\begin{equation}
\sigma = \pi d^{2}.
\end{equation}
Then consider the molecule traveling a distance $\langle v\rangle t$
(using the average velocity) will have a mean free path
\begin{equation}
\lambda = \frac{\begin{pmatrix}\mbox{distance}\\ \mbox{traveled}
  \end{pmatrix}}{\begin{pmatrix}\mbox{volume of}\\ \mbox{interaction}
  \end{pmatrix}\begin{pmatrix}\mbox{number of}\\ \mbox{molecules
      per}\\ \mbox{unit volume}
\end{pmatrix}}
= \frac{(\langle v\rangle t)}{(\pi d^{2}\cdot \langle v\rangle t)(n)} = \frac{1}{\pi d^{2}n}.
\end{equation}
This expression is a coarse approximation since we used the average
molecular velocity, and we implicitly assumed that the molecule hits
stationary targets. The frequency of collisions depends on the average
relative velocity of randomly moving molecules. This will change the
volume of interaction by a factor of $\sqrt{2}\approx1.41$.

The relative velocity between two molecules
\begin{equation}
\langle\vec{v}_{\text{rel}}^{2}\rangle=\langle(\vec{v}_{1}-\vec{v}_{2})^{2}\rangle=\langle\vec{v}_{1}^{2}+\vec{v}_{2}^{2}-2\vec{v}_{1}\cdot\vec{v}_{2}\rangle.
\end{equation}
But in equilibrium, $\vec{v}_{1}$ and $\vec{v}_{2}$ are random and
uncorrelated, so
\begin{equation}
\langle\vec{v}_{1}\cdot\vec{v}_{2}\rangle=0,
\end{equation}
and the relative speed is
\begin{equation}
v_{\text{rel}}=\sqrt{\langle\vec{v}_{\text{rel}}^{2}\rangle}=\sqrt{\langle\vec{v}_{1}^{2}+\vec{v}_{2}^{2}\rangle}=\sqrt{2}v.
\end{equation}
Then the volume of interaction is
\begin{equation}
\begin{pmatrix}\mbox{volume of}\\ \mbox{interaction}
  \end{pmatrix}=(\pi d^{2})(\sqrt{2}vt),
\end{equation}
which changes the mean-free path to
\begin{equation}
\lambda=\frac{1}{\sqrt{2}\pi d^{2}n}.
\end{equation}
\end{node}

\begin{xca}
(I thought this would be a good exercise---I do not have access to the
  class's problem sheets.)
Recall, in an ideal gas, we have:
\begin{equation}
pV = n_{\text{moles}}RT,
\end{equation}
and then Avogadro's number $N_{A}$ and the number of moles
$n_{\text{moles}}$ of substance gives us the number density:
\begin{equation}
n = \frac{n_{\text{moles}}N_{A}}{V}.
\end{equation}
Show that the mean free path for an ideal gas is then
\begin{equation}
\lambda = \frac{RT}{\sqrt{2}\pi d^{2}N_{A}p}.
\end{equation}
\end{xca}

\begin{definition}
The \define{Knudsen Number} $\Kn$ is defined to be the dimensionless
quantity $\Kn := \lambda/L$ where $\lambda$ is the mean free path and
$L$ is the length scale of the problem.
\end{definition}

\begin{node}
If $\Kn\ll1$, then we can use the continuum approach. (Usually,
$\Kn<0.1$ is the heuristic used.)

If $\Kn\geq0.1$, then the molecular/microscopic approach should be preferred.
\end{node}

\begin{remark}
In some older texts on gas dynamics, they demand even more: they
insist that $\Kn<0.001$ or so. Just be aware that the heuristic used
depends on the physical situation under consideration.
\end{remark}

\begin{example}
For air, let us consider when $L>10\lambda_{\text{air}}$. We
find that
\begin{subequations}
  \begin{align}
10\lambda_{\text{air}} &= \frac{10}{\sqrt{2}\pi d^{2}n}\\
&=\frac{10}{\sqrt{2}\pi d^{2}n_{0}}\left(\frac{n_{0}}{n}\right)\\
&=\left(\frac{n_{0}}{n}\right)(6.1\times10^{-7}~\mathrm{m}),
  \end{align}
\end{subequations}
or
\begin{equation}
\frac{L}{d}=1654\left(\frac{n_{0}}{n}\right).
\end{equation}
\end{example}

\subsection*{Meaningful statistics}

\begin{node}
  If we use the continuum approach, then (for example) mass density is
  \begin{equation}
\rho(\vec{x},t)=\frac{M}{V}
  \end{equation}
  as $V\to 0$ (where $M$ is the mass contained in the ``test region'' $V$).
\end{node}

\begin{node}
If we look at the probability density $P(N)$ describing the
probability of having $N$ molecules in volume $V$, then it is a
Poisson-distributed random variable:
\begin{equation}
P(N) = \frac{(nV)^{N}\exp(-nV)}{N!}.
\end{equation}
Its expected value is $nV$ and its variance is $nV$. Using the
``plus-or-minus $3\sigma$'' heuristic for describing ``most'' of the
molecules, we expect
\begin{equation}
N\sim nV\pm 3\sqrt{nV}.
\end{equation}
However, if $nV-3\sqrt{nV}\leq0$, then things go awry. This only
happens when $nV$ is ``small'' (say, $nV=1$ or smaller). For $nV\gg1$,
we'll be safe. So we have the last criterion be:
\begin{equation}
nV\gg1.
\end{equation}
We will continue discussing this heuristic next time.
\end{node}

\begin{xca}
Is the Poisson distribution applicable to this situation or not?
\end{xca}

\begin{remark}
The reference for the first and second lecture seems to be based
heavily on Bird~\cite[Ch~1]{bird1976molecular}.
\end{remark}

\endinput